\section{Welche Aufgabe habe ich im Praktikum}

Zu Beginn vom Praktikum war klar, dass ich nicht einfach nur irgendwie an einem Projekt arbeiten will, sondern mir eine saubere technische Basis aufbauen muss, die mir später das Arbeiten erleichtert. Dazu gehört eine gute Codeversionierung, technisches Verständnis der genutzten Technologien und ein stabiler GitHub-Workflow für die Dokumentation, die in LaTeX geschrieben wird, und die drei Bash-Skripte, die ich später für STEMgraph brauche, um automatisch mit Challenges und deren Abhängigkeiten arbeiten zu können.

Die Aufgaben in meinem Praktikum haben sich dabei in mehrere Bereiche aufgeteilt:

\begin{itemize}
\item STEMgraph.
\item Backend.
\item Frontend.
\item Security.
\item Deployment.
\end{itemize}

Das Ziel von diesem ersten abgeschlossenen Teil war für mich, ein LaTeX-Projektbericht mit Makefile aufzusetzen, das Repository mit einer sinnvollen \texttt{.gitignore} sauber zu halten, einen GitHub-Workflow zu bauen und Bash-Skripte zu schreiben, die STEMgraph-Challenges aus GitHub holen, rekursiv deren Abhängigkeiten auflösen und Teilpfade zwischen zwei Challenges berechnen.

Damit habe ich eine Grundlage geschaffen, auf der ich später das eigentliche STEMgraph-Frontend aufbaue und die Graphdaten exportieren und einbinden kann.

\subsection{STEMgraph}

STEMgraph ist das übergeordnete Projekt, in dem ich arbeite. Es geht dabei darum, Challenges, Abhängigkeiten und Lernpfade strukturiert abzubilden. Für mich war am Anfang wichtig zu verstehen, wie die Repositories aufgebaut sind und wie die Inhalte zusammenhängen.

\subsection{Backend}

Das Backend stellt die Daten und Funktionen bereit, die später vom Frontend genutzt werden. Dazu gehören die Verarbeitung der Challenge-Daten und die Verbindung zu den Graphinformationen. Für mich war dabei wichtig, die Datenflüsse im Gesamtsystem nachzuvollziehen.

\subsection{Frontend}

Beim Frontend geht es darum, die vorhandenen Inhalte verständlich darzustellen. Ich habe mir dazu Gedanken gemacht, wie man Lernpfade, Abhängigkeiten und Metadaten so anzeigt, dass man sich gut darin orientieren kann.

\subsection{Security}

Zum Bereich Security gehört vor allem der sichere Umgang mit Zugriffen, Rollen und geschützten Inhalten. Das ist relevant, weil nicht jede Information für jeden Nutzer gleich sichtbar sein soll.

\subsection{Deployment}

Beim Deployment geht es darum, wie die einzelnen Komponenten gebaut und bereitgestellt werden. Dafür sind Build-Prozesse, Releases und automatisierte Abläufe wichtig.